% Reflection positivity for the coupled slice (paper 12 of the O lane).
% Lane: O (operator bridge), module YangMills/OS/SpatialReflection.lean.
% REGIME: tricotomy; NO scores, NO desk language, NO commit chronology beyond
% the reproducibility freeze.
\documentclass[11pt]{article}
\usepackage[margin=1.1in]{geometry}
\usepackage{amsmath,amssymb,amsthm}
\usepackage{booktabs,longtable,array}
\usepackage[colorlinks=true,linkcolor=blue,urlcolor=blue,citecolor=blue]{hyperref}

\newtheorem{theorem}{Theorem}[section]
\newtheorem{lemma}[theorem]{Lemma}
\newtheorem{proposition}[theorem]{Proposition}
\newtheorem{corollary}[theorem]{Corollary}
\newtheorem{remark}[theorem]{Remark}
\theoremstyle{definition}
\newtheorem{definition}[theorem]{Definition}

\emergencystretch=3em
\tolerance=2000

\newcommand{\anchor}{c6a992a5d2abe453396989f40123487134f19e0c}
\newcommand{\osline}[3]{%
  \href{https://github.com/lluiseriksson/THE-ERIKSSON-PROGRAMME/blob/\anchor/YangMills/OS/#1.lean\#L#2}{\nolinkurl{#3}}}
\newcommand{\repofile}[2]{%
  \href{https://github.com/lluiseriksson/THE-ERIKSSON-PROGRAMME/blob/\anchor/#1}{\nolinkurl{#2}}}
\newcommand{\R}{\mathbb{R}}
\newcommand{\Z}{\mathbb{Z}}
\newcommand{\SU}{\mathrm{SU}}

\title{The Weight That Could Not Break It:\\
Machine-Checked Reflection Positivity\\
for the Coupled $\Z_2$ Slice\\
\large two reflections, two hypotheses, and one of them sharp}
\author{Lluis Eriksson}
\date{}

\begin{document}
\maketitle

\begin{abstract}
Eleven papers in this lane end with the same sentence: \emph{reflection
positivity is untouched}.  It is touched here, and the reason it is reachable is
the exact reason uniformity in the extent was not.

Everything this lane has proved about rates dies when the spatial source weight
is switched on: the modulus $\mathrm{specRatio}(L)$ runs to $1$ outside the
disordered region \cite{ErikssonSpectral}, and the extent-independent rate of
the companion paper holds only at constant weight \cite{ErikssonUniform}.
Reflection positivity is different in kind: conjugation by $\sqrt{w}$ is a
\emph{congruence}, and a congruence cannot change the sign of a quadratic form.
So the same weight that destroys uniformity leaves positivity exactly where it
was.  \textbf{This is the first statement in this lane that holds for the
coupled kernel.}

\textbf{Two reflections, two hypotheses.}  The reflected two-point sum is
$\langle v, K^{N} v\rangle$ with $v$ the dressed observable, and the two cases
have genuinely different content.  Through a \emph{site} ($N$ even) the sum is a
square, hence non-negative for \emph{every} $\beta$ --- negative coupling
included --- with nothing used but symmetry of the kernel.  Through a
\emph{bond} ($N$ odd) it is $\langle K^{m}v, K K^{m} v\rangle$, which needs $K$
itself positive semidefinite; that holds exactly for $\beta\ge0$, by an
induction on the extent.

\textbf{And the second hypothesis is active.}  At $\beta<0$ and odd separation
the single-site sign observable gives the reflected sum in closed form,
$2(e^{\beta}-e^{-\beta})^{N}$, which is negative.  The boundary $\beta=0$ is
therefore sharp and exhibited, not inferred.

\textbf{What this is not.}  Reflection positivity is one Osterwalder--Schrader
axiom.  \emph{No reconstruction is carried out}: the physical Hilbert space as
the quotient of the past algebra by the null space of this form is not built,
and nothing here is a statement about it.  The observables are the endpoint ones
the bridge module supplies \cite{ErikssonGibbs}; general past-algebra
observables are not treated.  Nothing here concerns uniformity in the extent,
$\SU(N)$, the continuum limit, or the Yang--Mills mass gap \cite{Clay}.

\textbf{On the pre-registration.}  Three gates were committed before any Lean was
written, and \S\ref{sec:gates} reports what happened to them --- including one
that failed because of a design error of ours, and the redesign that replaced
it.
\end{abstract}

\section{Why this one survives the weight}\label{sec:why}

The lane's obstruction has a shape.  Every rate statement it has produced is a
statement about $\mathrm{specGap}$ or $\mathrm{specRatio}$, and those are
spectral data: conjugating the kernel by $\sqrt{w}$ moves eigenvalues around,
and the measured evidence is that with a ring weight it moves them all the way
to degeneracy outside the disordered region.

Positivity is not spectral data in that sense.  It is the \emph{sign} of a
quadratic form, and
\[
  \langle u, \sqrt{w}\,T\,\sqrt{w}\, u\rangle
  \;=\; \langle \sqrt{w}\,u,\; T\, \sqrt{w}\,u\rangle
\]
is the same form read on a relabelled vector.  Sylvester's law is the general
statement; here it is one line, and it is the whole reason the coupled case
costs nothing once the decoupled one is done.

\section{Through a site: free, at every coupling}\label{sec:site}

\begin{theorem}[site reflection]\label{thm:site}
Let $K$ be any symmetric kernel.  For every $m$ and every observable $u$,
\[
  \sum_i \bigl(K^{2m} u\bigr)(i)\, u(i) \;\ge\; 0 .
\]
\end{theorem}

\osline{SpatialReflection}{73}{sum\_iterate\_even\_nonneg}.
Move half the iterates across the pairing --- which is what symmetry buys --- and
the sum becomes $\sum_i (K^{m}u)(i)^2$.  No positivity of $K$ is used, and no
hypothesis on $\beta$: this holds at negative coupling, where the bond
reflection below fails.

\section{Through a bond: exactly $\beta\ge0$}\label{sec:bond}

\begin{theorem}[the decoupled kernel is positive semidefinite]\label{thm:psd}
For $\beta\ge0$, every extent $L$ and every observable $u$,
$\sum_\sigma u(\sigma)\,(Ku)(\sigma) \ge 0$.
\end{theorem}

\osline{SpatialReflection}{115}{spatialKernel\_posSemidef}.
By induction on the extent, splitting an observable at the first site into its
even and odd parts.  Each part inherits positivity from the inductive
hypothesis, and the two are recombined with the two bond eigenvalues,
\[
  Z = e^{\beta}+e^{-\beta} > 0,
  \qquad
  D = e^{\beta}-e^{-\beta} \ge 0 \iff \beta \ge 0 ,
\]
which is precisely where the hypothesis enters.  It is the same induction the
companion uniformity paper runs \cite{ErikssonUniform}, with different
bookkeeping; the split lemma is shared and stated once
(\osline{SpatialReflection}{94}{act\_spatialKernel\_cons}).

\begin{theorem}[congruence, and the coupled kernel]\label{thm:congr}
If $K$ is positive semidefinite then so is $c\,K\,c$ for any function $c$.
Hence for $\beta\ge0$ and \emph{any} strictly positive source weight $w$, the
symmetrised kernel $\sqrt{w}\,K\,\sqrt{w}$ is positive semidefinite.
\end{theorem}

\osline{SpatialReflection}{177}{posSemidef\_congr} and
\osline{SpatialReflection}{193}{symWeighted\_posSemidef}.
This is the step \S\ref{sec:why} advertised, and it costs one reindexing.

\begin{theorem}[bond reflection]\label{thm:bond}
For $\beta\ge0$, every $m$ and every observable,
$\sum_i (K^{2m+1}u)(i)\,u(i) \ge 0$.
\end{theorem}

\osline{SpatialReflection}{207}{sum\_iterate\_odd\_nonneg}.

\section{The same statements, about the measure}\label{sec:measure}

The bridge of the companion paper \cite{ErikssonGibbs} identifies the reflected
two-point sum of the Gibbs measure with $\langle v, K^{N} v\rangle$, so both
theorems transfer verbatim.

\begin{theorem}[reflection positivity in the measure]\label{thm:rp}
For any strictly positive source weight, at every extent:
\begin{enumerate}
\item at even separation the reflected two-point sum is non-negative for
\emph{every} $\beta$
(\osline{SpatialReflection}{234}{gibbsPathSum\_nonneg\_even});
\item for $\beta\ge0$ it is non-negative at \emph{every} separation
(\osline{SpatialReflection}{243}{gibbsPathSum\_nonneg}).
\end{enumerate}
\end{theorem}

The statement is about the \emph{unnormalised} sum, which is the object
reflection positivity is about; the partition function is positive, so the
normalised expectation has the same sign, but that is a remark and not the
theorem.

\section{Sharpness: the hypothesis is active}\label{sec:sharp}

A hypothesis that is never in force is not worth stating.

\begin{theorem}[the witness, in closed form]\label{thm:witness}
At one site and constant weight, the single-site sign observable gives
\[
  \sum_X A(X_0)A(X_N)\,\text{weight}(X)
  \;=\; 2\,(e^{\beta}-e^{-\beta})^{N} ,
\]
exactly.  For $\beta<0$ and $N$ odd this is negative.
\end{theorem}

\osline{SpatialReflection}{275}{gibbsPathSum\_siteSign\_eq} and
\osline{SpatialReflection}{297}{gibbsPathSum\_neg\_of\_neg\_beta}.
The observable is the one the extent paper already built
\cite{ErikssonExtent}; it is an eigenvector of the decoupled kernel with
eigenvalue $D$, so the iterate is $D^{N}$ times it, and the norm counts the two
configurations.

\section{The three gates, and the one that failed}\label{sec:gates}

The previous campaign in this lane committed a judge whose single joint gate
read \textsc{fail} and was overridden; that is recorded as a process violation
in the ledger and is not undoable \cite{ErikssonUniform}.  The remedy applied
here was structural: separate gates, each licensing one theorem.

\begin{itemize}
\item \textbf{Gate A} (site reflection, every $\beta$): \textsc{pass}.
\item \textbf{Gate B} (bond reflection): \textsc{fail} --- and the reason is that
we made the same design error one level down.  Gate B bundled the $\beta\ge0$
theorem with its sharpness witness behind ``pass iff both''.  The first half
passed everywhere; the second missed one cell.
\end{itemize}

\repofile{scripts/judge\_spatial\_reflection.py}{scripts/judge\_spatial\_reflection.py}.
The autopsy measured why: at that cell the violating direction carries
$|D|^{N}/Z^{N} \approx 10^{-3}$ of the weight, so a uniformly random observable
lands in the violating cone with probability $\approx 2\times10^{-2}$, and sixty
draws miss it about a third of the time.  \textbf{An existence claim was being
tested by sampling.}
\repofile{scripts/judge\_spatial\_reflection\_autopsy.py}{scripts/judge\_spatial\_reflection\_autopsy.py}.

Gate B stays failed.  It was replaced by two new gates, committed before being
run and declared as a redesign rather than a reinterpretation: B1 keeps the
$\beta\ge0$ test, and B2 stops sampling and predicts a \emph{number} --- the
closed form of Theorem~\ref{thm:witness}, to $10^{-12}$ --- which is strictly
harder than asking some draw to go negative.  Both pass.
\repofile{scripts/judge\_spatial\_reflection\_v2.py}{scripts/judge\_spatial\_reflection\_v2.py}.

\section{Scope}\label{sec:scope}

\paragraph{One axiom, not a reconstruction.}  Reflection positivity is one of the
Osterwalder--Schrader axioms.  Building the physical Hilbert space --- the
quotient of the past algebra by the null space of this form, with the transfer
operator descending to a self-adjoint contraction --- is not done here and is
not claimed.

\paragraph{Endpoint observables only.}  The bridge module supplies correlations
of observables at the two ends of the path.  Reflection positivity for a general
observable of the whole past half-chain is not treated; the theorems above are
about the endpoint form.

\paragraph{Nothing about uniformity.}  This paper says nothing about how any
rate behaves as the extent grows.  The coupled kernel's uniformity is still the
lane's open analytic question.

\paragraph{No claim about the physical target.}  Nothing here concerns $\SU(N)$,
the continuum limit, or the Clay problem \cite{Clay}.  The distance to that
problem is unchanged.

\section{Reproducibility}\label{sec:repro}

All artifacts are at source checkpoint \texttt{\anchor} on branch \texttt{main}.
Every hyperlink was resolved against that commit, and checked to point at an
actual declaration line, before compilation.

\begin{center}
\begin{tabular}{ll}
\toprule
Quantity & Value \\
\midrule
Full core build & 8432 jobs, success \\
Oracle commands, and answers & 2708 \\
\quad distinct declarations & 2705 \\
\quad commands listed twice & 3 \\
\quad distinct, with axiom dependencies & 2682 \\
\quad distinct, axiom-free & 23 \\
Declarations in this module & 11 \\
Distinct axioms across all outputs & \texttt{propext}, \texttt{Quot.sound}, \texttt{Classical.choice} \\
Occurrences of \texttt{sorryAx} & 0 \\
Project axioms & 0 \\
Errors & 0 \\
\bottomrule
\end{tabular}
\end{center}

The module is
\repofile{YangMills/OS/SpatialReflection.lean}{YangMills/OS/SpatialReflection.lean};
verdicts are in
\repofile{docs/VERIFICATION-LEDGER.md}{docs/VERIFICATION-LEDGER.md}.

\begin{thebibliography}{9}

\bibitem{ErikssonExtent}
L.~Eriksson,
\emph{Where the Elementary Reconstruction Stops: a Machine-Checked Obstruction
at the Coupled $\Z_2$ Slice}, viXra:2607.0075, 2026.

\bibitem{ErikssonGibbs}
L.~Eriksson,
\emph{The Measure the Spectral Results Were About}, viXra, 2026.

\bibitem{ErikssonSpectral}
L.~Eriksson,
\emph{The Modulus: a Machine-Checked Operator Bound on the Fluctuation Sector of
the Coupled $\Z_2$ Slice}, viXra, 2026.

\bibitem{ErikssonUniform}
L.~Eriksson,
\emph{The Rate Without the Extent: a Machine-Checked Extent-Independent Spectral
Rate for the Decoupled $\Z_2$ Slice}, viXra, 2026.

\bibitem{OS}
K.~Osterwalder and R.~Schrader,
\emph{Axioms for Euclidean Green's functions}, Comm.\ Math.\ Phys.\
\textbf{31} (1973), 83--112.

\bibitem{OsterwalderSeiler}
K.~Osterwalder and E.~Seiler,
\emph{Gauge field theories on a lattice}, Ann.\ Physics \textbf{110} (1978),
440--471.

\bibitem{GlimmJaffe}
J.~Glimm and A.~Jaffe,
\emph{Quantum Physics: a Functional Integral Point of View}, 2nd ed.,
Springer, 1987.

\bibitem{Clay}
A.~Jaffe and E.~Witten,
\emph{Quantum Yang--Mills theory}, Clay Mathematics Institute Millennium Prize
problem description, 2000.

\end{thebibliography}

\end{document}
